% GENERATED FILE. Edit canonical lessons, then run npm run generate:books.
% canonical-source-hash: fnv1a64:255cd08b37b1c306
% canonical-lessons: GU-C489-kusal, GU-C489-sata, GU-C489-dua, GU-C489-sambharvu, GU-C489-hamdardi

\chapter{Asking After Someone}
\label{ch:gu-asking-after-someone}
\hlchaptermodality{489}

\begin{chapteropening}
\textbf{By the end of this chapter:} \emph{I can ask after someone's health and wish them well.}

Complete the last lesson of chapter 489: I can ask after someone's health and wish them well.
\end{chapteropening}

\section[\texorpdfstring{\gu{કુશળ} (kuśaḷ)}{કુશળ (kuśaḷ)}]{\gu{કુશળ} (kuśaḷ) — faring well; safe}

\label{lesson:GU-C489-kusal}



\begin{quote}
Before the new one: say the Gujarati for to play with a child, then the Gujarati for to give someone a drink.
\end{quote}

\subsection*{You'll want to know: \gu{કુશળ}}
\textbf{\gu{કુશળ}} — \emph{kuśaḷ} — \textquotedblleft{}faring well; safe\textquotedblright{}.

\begin{quote}
\gu{ઘરમાં} \gu{બધા} \gu{કુશળ} \gu{છે} \gu{કે} \gu{નહીં}?
\end{quote}

\begin{quote}
\emph{gharmā̃ badhā kuśaḷ che ke nahī̃?}
\end{quote}

\begin{quote}
Is everyone at home well?
\end{quote}

\begin{cousinweb}
\textbf{\gu{કુશળ}} (\emph{kuśaḷ}) is a tatsama, Sanskrit \emph{kuśala}, \textquotedblleft{}right, proper\textquotedblright{} and \textquotedblleft{}well, healthy\textquotedblright{}, as Monier-Williams records it.
\end{cousinweb}

\begin{grammarlens}[title={What you've built}]
One more word for describing what you lived and felt.
\end{grammarlens}

\subsection*{Your turn}
\begin{itemize}
\raggedright
  \item \emph{Say it:} \emph{kuśaḷ}
  \item \emph{Say it:} \emph{kuśaḷ}, once more
  \item \emph{Say it:} \emph{gharmā̃ badhā kuśaḷ che ke nahī̃?} — Is everyone at home well?
  \item \emph{Recall:} say the Gujarati for to play with a child, then the Gujarati for to give someone a drink, then say \emph{kuśaḷ} again
\end{itemize}

\begin{tcolorbox}[breakable,colback=teal!4,colframe=teal!35!black,title={Before you move on}]
What does \gu{કુશળ} mean? (\textquotedblleft{}Faring well; safe\textquotedblright{}.) Say it once more.
\end{tcolorbox}
\section[\texorpdfstring{\gu{શાતા} (śātā)}{શાતા (śātā)}]{\gu{શાતા} (śātā) — ease; relief from pain}

\label{lesson:GU-C489-sata}



\begin{quote}
Before the new one: say the Gujarati for to give someone a drink, then the Gujarati for faring well.
\end{quote}

\subsection*{You'll want to know: \gu{શાતા}}
\textbf{\gu{શાતા}} — \emph{śātā} — \textquotedblleft{}ease; relief from pain\textquotedblright{}.

\begin{quote}
\gu{મીરાએ} \gu{ફોન} \gu{કરીને} \gu{દાદીની} \gu{શાતા} \gu{પૂછી}.
\end{quote}

\begin{quote}
\emph{Mīrāe phon karīne dādīnī śātā pūchī.}
\end{quote}

\begin{quote}
Mira phoned and asked how her grandmother was feeling.
\end{quote}

A feminine noun.

\begin{cousinweb}
\textbf{\gu{શાતા}} (\emph{śātā}) probably goes back to Sanskrit \emph{śāta}, \textquotedblleft{}joy, pleasure, happiness\textquotedblright{}, which Monier-Williams records from the old word lists; Gujarati uses it for ease after pain, and \textbf{\gu{શાતા} \gu{પૂછવી}} is to ask after someone who is ill.
\end{cousinweb}

\begin{grammarlens}[title={What you've built}]
One more word for describing what you lived and felt.
\end{grammarlens}

\subsection*{Your turn}
\begin{itemize}
\raggedright
  \item \emph{Say it:} \emph{śātā}
  \item \emph{Say it:} \emph{śātā}, once more
  \item \emph{Say it:} \emph{Mīrāe phon karīne dādīnī śātā pūchī.} — Mira phoned and asked how her grandmother was feeling.
  \item \emph{Recall:} say the Gujarati for to give someone a drink, then the Gujarati for faring well, then say \emph{śātā} again
\end{itemize}

\begin{tcolorbox}[breakable,colback=teal!4,colframe=teal!35!black,title={Before you move on}]
What does \gu{શાતા} mean? (\textquotedblleft{}Ease; relief from pain\textquotedblright{}.) Say it once more.
\end{tcolorbox}
\section[\texorpdfstring{\gu{દુઆ} (duā)}{દુઆ (duā)}]{\gu{દુઆ} (duā) — a prayer for someone; a blessing}

\label{lesson:GU-C489-dua}



\begin{quote}
Before the new one: say the Gujarati for faring well, then the Gujarati for relief from pain.
\end{quote}

\subsection*{You'll want to know: \gu{દુઆ}}
\textbf{\gu{દુઆ}} — \emph{duā} — \textquotedblleft{}a prayer for someone; a blessing\textquotedblright{}.

\begin{quote}
\gu{દાદીની} \gu{દુઆથી} \gu{મારો} \gu{તાવ} \gu{મટી} \gu{ગયો}.
\end{quote}

\begin{quote}
\emph{dādīnī duāthī māro tāv maṭī gayo.}
\end{quote}

\begin{quote}
Thanks to my grandmother's prayers my fever went away.
\end{quote}

A feminine noun.

\begin{cousinweb}
\textbf{\gu{દુઆ}} (\emph{duā}) is usually traced to Arabic \emph{duʿā}, \textquotedblleft{}a prayer, a call\textquotedblright{}, which reached Gujarati through Persian.
\end{cousinweb}

\begin{grammarlens}[title={What you've built}]
One more word for describing what you lived and felt.
\end{grammarlens}

\subsection*{Your turn}
\begin{itemize}
\raggedright
  \item \emph{Say it:} \emph{duā}
  \item \emph{Say it:} \emph{duā}, once more
  \item \emph{Say it:} \emph{dādīnī duāthī māro tāv maṭī gayo.} — Thanks to my grandmother's prayers my fever went away.
  \item \emph{Recall:} say the Gujarati for faring well, then the Gujarati for relief from pain, then say \emph{duā} again
\end{itemize}

\begin{tcolorbox}[breakable,colback=teal!4,colframe=teal!35!black,title={Before you move on}]
What does \gu{દુઆ} mean? (\textquotedblleft{}A prayer for someone; a blessing\textquotedblright{}.) Say it once more.
\end{tcolorbox}
\section[\texorpdfstring{\gu{સંભારવું} (sambhārvũ)}{સંભારવું (sambhārvũ)}]{\gu{સંભારવું} (sambhārvũ) — to recall; to remember (someone) fondly}

\label{lesson:GU-C489-sambharvu}



\begin{quote}
Before the new one: say the Gujarati for relief from pain, then the Gujarati for a prayer for someone.
\end{quote}

\subsection*{You'll want to know: \gu{સંભારવું}}
\textbf{\gu{સંભારવું}} — \emph{sambhārvũ} — \textquotedblleft{}to recall; to remember (someone) fondly\textquotedblright{}.

\begin{quote}
\gu{દાદા} \gu{રોજ} \gu{જૂના} \gu{મિત્રોને} \gu{સંભારે} \gu{છે}.
\end{quote}

\begin{quote}
\emph{dādā roj jūnā mitrone sambhāre che.}
\end{quote}

\begin{quote}
Every day my grandfather remembers his old friends.
\end{quote}

Past, with \gu{તેણે} (he or she …): \gu{તેણે} \gu{જૂના} \gu{દિવસો} \gu{સંભાર્યા}, \emph{teṇe jūnā divso sambhāryā}.

\begin{cousinweb}
\textbf{\gu{સંભારવું}} (\emph{sambhārvũ}) is inherited from Sanskrit \emph{saṃsmārayati}, \textquotedblleft{}causes to remember\textquotedblright{}; CDIAL lists the Gujarati verb, \textquotedblleft{}to recollect\textquotedblright{} (CDIAL 13059).
\end{cousinweb}

\begin{grammarlens}[title={What you've built}]
One more word for describing what you lived and felt.
\end{grammarlens}

\subsection*{Your turn}
\begin{itemize}
\raggedright
  \item \emph{Say it:} \emph{sambhārvũ}
  \item \emph{Say it:} \emph{sambhārvũ}, once more
  \item \emph{Say it:} \emph{dādā roj jūnā mitrone sambhāre che.} — Every day my grandfather remembers his old friends.
  \item \emph{Recall:} say the Gujarati for relief from pain, then the Gujarati for a prayer for someone, then say \emph{sambhārvũ} again
\end{itemize}

\begin{tcolorbox}[breakable,colback=teal!4,colframe=teal!35!black,title={Before you move on}]
What does \gu{સંભારવું} mean? (\textquotedblleft{}To recall; to remember (someone) fondly\textquotedblright{}.) Say it once more.
\end{tcolorbox}
\section[\texorpdfstring{\gu{હમદર્દી} (hamdardī)}{હમદર્દી (hamdardī)}]{\gu{હમદર્દી} (hamdardī) — sympathy}

\label{lesson:GU-C489-hamdardi}



\begin{quote}
Before the new one: say the Gujarati for a prayer for someone, then the Gujarati for to recall.
\end{quote}

\subsection*{You'll want to know: \gu{હમદર્દી}}
\textbf{\gu{હમદર્દી}} — \emph{hamdardī} — \textquotedblleft{}sympathy\textquotedblright{}.

\begin{quote}
\gu{પડોશીઓએ} \gu{અમારા} \gu{દુઃખમાં} \gu{હમદર્દી} \gu{બતાવી}.
\end{quote}

\begin{quote}
\emph{paḍośīoe amārā duḥkhmā̃ hamdardī batāvī.}
\end{quote}

\begin{quote}
The neighbours showed sympathy in our sorrow.
\end{quote}

A feminine noun.

\begin{cousinweb}
\textbf{\gu{હમદર્દી}} (\emph{hamdardī}) is usually traced to Persian \emph{hamdardī}, \textquotedblleft{}sharing pain\textquotedblright{}, built from \emph{ham}, \textquotedblleft{}together\textquotedblright{}, and \emph{dard}, \textquotedblleft{}pain\textquotedblright{}.
\end{cousinweb}

\begin{grammarlens}[title={What you've built}]
One more word for describing what you lived and felt.
\end{grammarlens}

\subsection*{Your turn}
\begin{itemize}
\raggedright
  \item \emph{Say it:} \emph{hamdardī}
  \item \emph{Say it:} \emph{hamdardī}, once more
  \item \emph{Say it:} \emph{paḍośīoe amārā duḥkhmā̃ hamdardī batāvī.} — The neighbours showed sympathy in our sorrow.
  \item \emph{Recall:} say the Gujarati for a prayer for someone, then the Gujarati for to recall, then say \emph{hamdardī} again
\end{itemize}

\begin{tcolorbox}[breakable,colback=teal!4,colframe=teal!35!black,title={Before you move on}]
What does \gu{હમદર્દી} mean? (\textquotedblleft{}Sympathy\textquotedblright{}.) Say it once more.
\end{tcolorbox}
